\documentclass[11pt,a4paper]{article}

\usepackage[utf8]{inputenc}
\usepackage[T1]{fontenc}
\usepackage[french]{babel}
\usepackage{lmodern}
\usepackage{microtype}
\usepackage[margin=2.3cm]{geometry}
\usepackage{booktabs}
\usepackage{tabularx}
\usepackage{array}
\usepackage{xcolor}
\usepackage{pifont}
\usepackage{enumitem}
\usepackage{float}
\usepackage[round,authoryear]{natbib}
\usepackage[colorlinks=true,linkcolor=teal!70!black,citecolor=teal!70!black,urlcolor=teal!70!black]{hyperref}

\newcommand{\oui}{\textcolor{teal!70!black}{\ding{51}}}
\newcommand{\non}{\textcolor{red!60!black}{\ding{55}}}
\newcommand{\partiel}{$\sim$}
\newcolumntype{Y}{>{\raggedright\arraybackslash}X}
\newcolumntype{C}[1]{>{\centering\arraybackslash}p{#1}}
\setlist{itemsep=2pt,topsep=4pt}

\title{\textbf{Livrable L1 -- État de l'art}\\[4pt]
\large Allocation adaptative des ressources de calcul et de recherche documentaire\\ pour des systèmes LLM locaux, économes et fiables\\[2pt]
\normalsize\textit{Adaptive Compute and Retrieval Allocation for Energy-Efficient and Reliable Local LLM Systems}}
\author{Alae Lahbichi -- Étudiant en 5\textsuperscript{e} année IA/Data\\[2pt]
\small\href{mailto:a.lahbichi2004@gmail.com}{\texttt{a.lahbichi2004@gmail.com}}\\
\small Dépôt GitHub : \href{https://github.com/AlaeLahbichi/Adaptive-Compute-and-RAG-for-Efficient-Reliable-Local-LLMs}{\footnotesize\texttt{AlaeLahbichi/Adaptive-Compute-and-RAG-for-Efficient-Reliable-Local-LLMs}}}
\date{Version de travail v0.1 -- 17 octobre 2026}

\begin{document}
\maketitle

\begin{center}
\fcolorbox{teal!60!black}{teal!5}{\parbox{0.92\textwidth}{\small
\textbf{Statut du document.} Cette version a été rédigée à partir des métadonnées et des résumés
des articles recensés ; les lectures approfondies (passes 2 et 3) restent à faire. Les affirmations
sur le contenu des articles doivent être vérifiées dans les textes complets avant toute réutilisation
dans un article. Le positionnement proposé (section~\ref{sec:positionnement}) est une
\emph{hypothèse de nouveauté} à confirmer par la lecture complète des travaux les plus proches, et
non une nouveauté établie.}}
\end{center}

\begin{abstract}
\noindent Les systèmes fondés sur des grands modèles de langage (LLM) traitent généralement toutes les
requêtes avec la même configuration : un même modèle et, le cas échéant, une recherche documentaire (RAG)
identique. Cette uniformité gaspille des ressources sur les questions simples et peut dégrader la
fiabilité sur les questions difficiles ou non couvertes par les documents. Ce document recense
51 travaux, dont 46 retenus, répartis en cinq axes : le routage et les cascades de modèles, le RAG
adaptatif, l'estimation d'incertitude et l'abstention, la mesure du coût énergétique de l'inférence,
et les jeux d'évaluation. L'analyse montre que ces trois décisions -- \emph{quel modèle},
\emph{faut-il chercher et combien}, \emph{faut-il répondre} -- sont le plus souvent étudiées
séparément, et que le coût est majoritairement exprimé en prix d'API ou en nombre d'appels plutôt
qu'en latence, mémoire ou énergie mesurées sur du matériel local. Nous en déduisons un positionnement
possible : une politique de décision \emph{conjointe} et interprétable, évaluée sur machine locale.
Des travaux de 2026 réduisent toutefois cet écart, ce qui impose de délimiter précisément la
contribution.
\end{abstract}

\clearpage
\tableofcontents
\clearpage

%======================================================================
\section{Introduction}
%======================================================================

\paragraph{Contexte.} Exécuter un LLM localement présente deux avantages : les données restent sur la
machine et le coût est maîtrisé. Les ressources d'un ordinateur personnel sont en revanche limitées en
mémoire, en temps de calcul et en énergie. Or toutes les questions ne demandent pas le même effort :
certaines sont triviales pour un petit modèle, d'autres exigent un modèle plus puissant, des documents
externes, voire un refus de répondre lorsque les preuves manquent.

\paragraph{Problématique.} \emph{Comment allouer automatiquement, et de manière mesurable, les ressources
de calcul et de recherche documentaire à chaque requête, afin de réduire le coût d'un système LLM local
sans dégrader significativement la qualité, la factualité et la fiabilité des réponses ?}

\paragraph{Objectifs de cette revue.} Elle doit (i) recenser les approches existantes pour chacune des
décisions d'allocation, (ii) identifier les signaux utilisés pour décider, (iii) examiner comment le coût
et la fiabilité sont évalués, et (iv) en déduire ce qui distinguerait réellement la méthode envisagée.
Le tableau~\ref{tab:qr} rattache chaque axe aux questions de recherche du cahier des charges.

\begin{table}[h]
\centering\small
\begin{tabularx}{\textwidth}{l Y l}
\toprule
\textbf{QR} & \textbf{Question} & \textbf{Axes de la revue} \\
\midrule
QR1 & Peut-on estimer la difficulté d'une question avant de choisir le modèle ou de lancer le RAG ? & §\ref{sec:routage}, §\ref{sec:rag} \\
QR2 & Des signaux d'incertitude et de qualité des preuves permettent-ils de décider quand chercher davantage ou s'abstenir ? & §\ref{sec:rag}, §\ref{sec:incertitude} \\
QR3 & Quel compromis qualité / latence / mémoire / énergie par rapport à des stratégies fixes ? & §\ref{sec:energie} \\
QR4 & Les gains sont-ils stables selon les types de questions et de documents ? & §\ref{sec:benchmarks} \\
\bottomrule
\end{tabularx}
\caption{Correspondance entre les questions de recherche et les axes de la revue.}
\label{tab:qr}
\end{table}

\paragraph{Plan.} La section~\ref{sec:methode} décrit la méthode de recherche bibliographique, la
section~\ref{sec:notions} définit les notions de base, et les sections~\ref{sec:routage}
à~\ref{sec:benchmarks} présentent chaque axe. La section~\ref{sec:comparaison} propose un tableau
comparatif, puis la section~\ref{sec:positionnement} identifie le manque et formule la contribution
potentielle. La section~\ref{sec:limites} précise les limites de cette revue.

%======================================================================
\section{Méthode de recherche bibliographique}\label{sec:methode}
%======================================================================

\paragraph{Sources.} arXiv, ACL Anthology, OpenReview, Google Scholar, Semantic Scholar et les actes
(proceedings) de NeurIPS, ICLR, AAAI et MLSys. Les références sont gérées dans Zotero.

\paragraph{Requêtes.} Elles combinent les concepts du sujet et leurs synonymes : \textit{LLM routing},
\textit{model cascade}, \textit{adaptive retrieval}, \textit{adaptive RAG}, \textit{dynamic RAG},
\textit{uncertainty estimation}, \textit{selective prediction}, \textit{abstention},
\textit{energy-aware LLM inference}. Une recherche ciblée sur arXiv (6 requêtes croisées, menée le
17/10/2026, détaillée en annexe~\ref{ann:requetes}) a servi à repérer les travaux de 2025--2026 qui
combinent plusieurs de ces décisions.

\paragraph{Critères d'inclusion.} (i) La méthode adapte, requête par requête, le choix du modèle, la
récupération documentaire ou la décision de répondre ; (ii) elle s'accompagne d'une évaluation
expérimentale ; ou (iii) l'article fournit une mesure du coût d'inférence ou un jeu de données utile
au protocole.

\paragraph{Critères d'exclusion.} Optimisation d'infrastructure cloud, architectures matérielles,
articles de position sans expérience et systèmes multi-agents hors du cadre questions-réponses.

\paragraph{Processus de sélection.}
\begin{itemize}
\item 17 références initiales, issues du cahier des charges et d'une première recherche ;
\item 26 références ajoutées par recherche ciblée et par \textit{snowballing}, c'est-à-dire en
suivant les références citées et citantes ;
\item 46 résultats arXiv examinés sur titre lors de la recherche croisée, dont 8 retenus après lecture
du résumé ;
\item soit \textbf{51 références identifiées}, dont \textbf{5 exclues} (hors périmètre) et
\textbf{46 retenues}, classées par thème et par priorité de lecture (P1 à P3).
\end{itemize}

%======================================================================
\section{Notions de base}\label{sec:notions}
%======================================================================

\begin{description}[style=nextline,leftmargin=0.5cm]
\item[Routage (\textit{routing})] Choix, avant la génération, du modèle qui traitera une requête parmi
plusieurs candidats de tailles ou de coûts différents.
\item[Cascade] Variante séquentielle du routage : on interroge d'abord un modèle peu coûteux, puis on
escalade vers un modèle plus puissant si sa réponse semble peu fiable. Le coût du premier modèle est
toujours payé.
\item[RAG] Génération augmentée par récupération : des passages extraits d'une base documentaire sont
ajoutés au prompt. Le RAG \emph{adaptatif} décide si, quand et combien récupérer.
\item[Incertitude] Estimation de la probabilité qu'une réponse soit incorrecte : probabilités internes,
cohérence entre plusieurs réponses échantillonnées, ou confiance exprimée par le modèle.
\item[Abstention] Refus de répondre lorsque le risque d'erreur est jugé trop élevé. On l'évalue par le
compromis \emph{couverture / risque} : la part de questions auxquelles le système répond, et le taux
d'erreur parmi ces réponses.
\item[Coût] Selon les travaux : prix d'API, nombre d'appels au grand modèle, nombre de tokens, latence,
mémoire ou énergie mesurée. Ces mesures ne sont pas interchangeables.
\end{description}

%======================================================================
\section{Routage et cascades de modèles}\label{sec:routage}
%======================================================================

\paragraph{Cascades.} \citet{chen2023frugalgpt} distinguent trois leviers pour réduire le coût des LLM :
l'adaptation du prompt, l'approximation du LLM et la cascade. Leur instanciation, FrugalGPT, apprend
quelles combinaisons de modèles utiliser selon la requête, et rapporte jusqu'à 98\,\% d'économie à
performance égale à celle du meilleur modèle individuel. AutoMix \citep{aggarwal2023automix} escalade
vers un modèle plus grand selon une auto-vérification \textit{few-shot} de la réponse du petit modèle ;
comme cette auto-vérification est bruitée, la décision est confiée à un routeur fondé sur un processus
de décision markovien partiellement observable (POMDP), avec plus de 50\,\% d'économie à performance
comparable. \citet{yue2023cascades} utilisent comme signal de difficulté la \emph{cohérence} entre
plusieurs réponses du petit modèle et atteignent la performance du modèle fort pour 40\,\% de son coût.

\paragraph{Routage avant génération.} Hybrid LLM \citep{ding2024hybrid} prédit la difficulté de la
requête pour l'attribuer à un petit modèle (déployable en périphérie, dit \textit{edge}) ou à un grand
modèle. Un seuil de qualité réglable à l'inférence permet d'éviter jusqu'à 40\,\% des appels au grand
modèle sans baisse de qualité. RouteLLM \citep{ong2024routellm} entraîne des routeurs fort/faible à
partir de données de préférences humaines. Le coût est réduit de plus de moitié dans certains cas, et
les routeurs restent efficaces lorsque la paire de modèles change. D'autres formulations existent : le
bandit manchot \citep{nguyen2024metallm} ou le routage vers le modèle expert guidé par un modèle de
récompense \citep{lu2023zooter}. RouterBench \citep{hu2024routerbench} propose un cadre d'évaluation
standardisé et un corpus de plus de 405\,000 résultats d'inférence.

\paragraph{Travaux récents.} \citet{fu2026inflation} montrent qu'un routage fondé sur le coût d'un
appel unique sous-estime le coût réel d'un workflow qui relance les requêtes, jusqu'à un facteur
supérieur à 2 sur les tâches difficiles. Ils introduisent un signal de difficulté calculé
\emph{localement avant exécution} et sélectionnent le modèle selon le rapport entre précision attendue
et coût réel prédit. RACER \citep{hao2026racer} étend un routeur pour produire des \emph{ensembles} de
modèles, avec un seuil calibré qui autorise l'abstention et contrôle le risque d'erreur de routage
(garanties sans hypothèse sur la distribution des données). La revue de \citet{varangot2025survey}
synthétise le domaine et souligne deux limites : la standardisation insuffisante des expériences et la
faible prise en compte des coûts non financiers.

\paragraph{Bilan.} Ces travaux choisissent \emph{le modèle}, mais ne décident ni de la récupération
documentaire ni, à l'exception de RACER, de l'abstention. Le coût y est surtout exprimé en prix d'API
ou en nombre d'appels.

%======================================================================
\section{RAG adaptatif}\label{sec:rag}
%======================================================================

\paragraph{Quand récupérer ?} \citet{mallen2022trust} montrent que les LLM retiennent mal les
connaissances peu populaires. Récupérer \emph{uniquement lorsque c'est nécessaire} améliore alors à la
fois la précision et le coût. Les approches suivantes diffèrent par le signal de décision :
\begin{itemize}
\item la connaissance que le modèle a de lui-même, estimée par comparaison avec des questions déjà
traitées (SKR, \citealp{wang2023skr}) ;
\item la confiance des tokens générés, pour déclencher la récupération en cours de génération (FLARE,
\citealp{jiang2023flare}) ;
\item une estimation du besoin d'information en temps réel (DRAGIN, \citealp{su2024dragin}) ;
\item des tokens de réflexion appris par le modèle, ce qui suppose un entraînement spécifique
(Self-RAG, \citealp{asai2023selfrag}).
\end{itemize}

\paragraph{Selon la complexité de la question.} Adaptive-RAG \citep{jeong2024adaptiverag} entraîne un
petit classifieur qui prédit la complexité d'une requête pour choisir entre trois stratégies : pas de
récupération, récupération en une étape, ou récupération itérative. Le système gagne à la fois en
efficacité et en précision sur des jeux de questions simples et multi-étapes. C'est la référence la
plus directe pour QR1, et son code est public. HyPA-RAG \citep{kalra2024hyparag} applique la même
idée aux paramètres de récupération dans le domaine juridique.

\paragraph{Selon la qualité des preuves.} Corrective RAG \citep{yan2024crag} évalue la qualité des
documents récupérés avec un évaluateur léger, puis déclenche une action selon le résultat : utiliser les
documents, les compléter par une recherche web, ou les filtrer. Tiny-Critic RAG \citep{wu2026tinycritic}
confie ce rôle de contrôle à un petit modèle adapté par LoRA (adaptation à faible rang), avec une latence
réduite d'un ordre de grandeur par rapport à un grand évaluateur.

\paragraph{Incertitude contre pipelines complexes.} \citet{moskvoretskii2025adaptive} comparent
35 méthodes (8 méthodes de récupération adaptative et 27 estimateurs d'incertitude) sur 6 jeux de
données et 10 métriques. Les estimateurs d'incertitude simples font souvent mieux que les pipelines
complexes en efficacité, pour une qualité comparable. Ce résultat justifie de privilégier des
\emph{signaux simples et interprétables}. Self-Routing RAG \citep{wu2025srrag} généralise la décision
binaire « récupérer ou non » en un choix de \emph{source de connaissance}, dans lequel le LLM lui-même
est une source. Avec des modèles de taille 7B, il réduit de 21 à 40\,\% le nombre de récupérations.

\paragraph{Routage entre paradigmes de RAG.} RAGRouter-Bench \citep{wang2026ragrouterbench} formule le
choix d'un paradigme de RAG comme un problème de routage dépendant de la paire requête--corpus. Il
évalue conjointement la qualité et la consommation de ressources avec plusieurs LLM, et conclut
qu'aucun paradigme n'est optimal partout. Sur ce benchmark, un simple classifieur TF-IDF + SVM simule
28\,\% d'économie de tokens \citep{bansal2026lightweight}. Pre-Route \citep{chen2026preroute} choisit de
façon analogue entre RAG et contexte long.

\paragraph{Réduire le contexte.} LeanContext \citep{arefeen2024leancontext} réduit le contexte récupéré
aux phrases clés : le coût baisse de 37 à 68\,\% pour une perte de 1,4 à 2,7\,\% de ROUGE-1. Ce
compromis chiffré est du même ordre que le critère de succès du projet. Clustered RAG
\citep{akesson2024clustered} et LogicRAG \citep{chen2026logicrag} réduisent eux aussi les tokens de
contexte.

\paragraph{Bilan.} Le RAG adaptatif décide \emph{si et comment récupérer}, généralement avec \emph{un
seul modèle générateur}. L'abstention y est rarement une sortie explicite, et le coût y est le plus
souvent mesuré en nombre de récupérations ou d'appels.

%======================================================================
\section{Estimation d'incertitude et abstention}\label{sec:incertitude}
%======================================================================

\paragraph{Estimer l'incertitude.} Trois familles de signaux se distinguent.
\begin{itemize}
\item \emph{Probabilités internes.} \citet{kadavath2022know} montrent que les grands modèles sont
souvent bien calibrés lorsqu'on leur demande d'évaluer leurs propres réponses (P(True)). La probabilité
« je connais la réponse », P(IK), augmente en présence de sources pertinentes, ce qui en fait un signal
candidat pour décider de récupérer.
\item \emph{Cohérence entre plusieurs échantillons.} L'entropie sémantique \citep{kuhn2023semantic}
regroupe les réponses échantillonnées qui ont le même sens ; SelfCheckGPT \citep{manakul2023selfcheckgpt}
détecte les hallucinations par contradiction entre échantillons. Ces signaux sont efficaces mais
exigent \emph{plusieurs générations} par question.
\item \emph{Confiance verbalisée.} Le modèle exprime sa propre confiance en mots ; elle n'est calibrée
qu'après un apprentissage dédié \citep{lin2022verbal}.
\end{itemize}

\paragraph{S'abstenir.} \citet{kamath2020selective} formalisent la réponse sélective sous changement de
domaine. Les probabilités \textit{softmax} sont trop confiantes hors domaine, alors qu'un calibrateur
appris améliore nettement la couverture à précision fixée. La revue de \citet{wen2024abstention}
organise les méthodes, benchmarks et métriques d'abstention selon trois perspectives : la requête, le
modèle et les valeurs humaines.

\paragraph{Bilan.} Le coût de l'estimation d'incertitude elle-même est rarement intégré à la décision.
Pourtant, échantillonner plusieurs réponses avec un petit modèle peut coûter plus cher qu'un appel
direct à un modèle plus grand. Ce point concerne directement QR2 et QR3.

%======================================================================
\section{Coût, efficacité et énergie de l'inférence}\label{sec:energie}
%======================================================================

\citet{samsi2023watts} mesurent l'énergie de l'inférence de LLaMA selon la taille du modèle, sur des
GPU de centre de données (V100, A100). \citet{luccioni2022bloom} estiment l'empreinte carbone de BLOOM
sur l'ensemble de son cycle de vie, et \citet{luccioni2024power} comparent le coût énergétique de
l'inférence de nombreux modèles et tâches : les modèles génératifs généralistes s'y révèlent beaucoup
plus coûteux que les modèles spécialisés.

\paragraph{Bilan.} Ces travaux \emph{mesurent} l'énergie mais ne l'\emph{optimisent} pas par des
décisions prises requête par requête. Ils portent en outre sur du matériel de centre de données.
Inversement, les travaux de routage optimisent un coût qui n'est presque jamais énergétique. Le cahier
des charges impose une précaution que la littérature confirme : ne pas déduire une énergie exacte du
seul temps CPU, et rapporter séparément les indicateurs indirects (latence, tokens, mémoire).

%======================================================================
\section{Jeux de données et protocoles d'évaluation}\label{sec:benchmarks}
%======================================================================

Les jeux de données recensés couvrent les catégories de questions demandées par le cahier des charges
(tableau~\ref{tab:data}).

\begin{table}[h]
\centering\small
\begin{tabularx}{\textwidth}{l Y}
\toprule
\textbf{Catégorie} & \textbf{Jeux de données} \\
\midrule
Simple & SQuAD (questions avec réponse) \citep{rajpurkar2018squad2} \\
Complexe / multi-étapes & HotpotQA, avec phrases justificatives \citep{yang2018hotpotqa} ; MuSiQue \citep{trivedi2022musique} \\
Non couverte par les documents & SQuAD 2.0 \citep{rajpurkar2018squad2} ; MuSiQue-Full (questions sans réponse) \\
Types variés, hallucinations pénalisées & CRAG Benchmark \citep{yang2024cragbench} \\
Routage entre paradigmes de RAG & RAGRouter-Bench \citep{wang2026ragrouterbench} \\
\bottomrule
\end{tabularx}
\caption{Jeux de données utilisables pour le benchmark du projet (Phase 2).}
\label{tab:data}
\end{table}

\paragraph{Leçons méthodologiques.} (i) \citet{li2026beyond} montrent que des signaux issus de la
récupération, pourtant corrélés à l'utilité d'une étape supplémentaire, n'améliorent pas de façon
fiable un routeur fondé sur la seule requête. Tout signal ajouté doit donc être comparé à un
\emph{contrôle « requête seule »} apparié. (ii) RouterBench et RAGRouter-Bench évaluent la qualité et le
coût conjointement, sous forme de compromis, plutôt qu'avec une seule métrique. (iii) Adaptive-RAG
s'évalue sur SQuAD, HotpotQA et MuSiQue : réutiliser ces jeux permet une comparaison directe.

%======================================================================
\section{Tableau comparatif}\label{sec:comparaison}
%======================================================================

Le tableau~\ref{tab:comparaison} situe les travaux les plus proches selon les trois décisions
d'allocation, le signal utilisé, la nature du coût évalué et le cadre d'exécution.

\begin{table}[H]
\centering\footnotesize
\setlength{\tabcolsep}{3.5pt}
\begin{tabularx}{\textwidth}{l C{1.25cm} C{1.25cm} C{1.35cm} Y Y C{0.95cm}}
\toprule
\textbf{Travail} & \textbf{Choix du modèle} & \textbf{Décision RAG} & \textbf{Abs\-ten\-tion} & \textbf{Signal de décision} & \textbf{Coût évalué} & \textbf{Local} \\
\midrule
FrugalGPT \citeyearpar{chen2023frugalgpt} & \oui & \non & \non & Scoreur appris (après génération) & Prix API & \non \\
Hybrid LLM \citeyearpar{ding2024hybrid} & \oui & \non & \non & Difficulté prédite (avant) & Appels au grand modèle & \partiel \\
RouteLLM \citeyearpar{ong2024routellm} & \oui & \non & \non & Préférences humaines & Prix / appels & \non \\
AutoMix \citeyearpar{aggarwal2023automix} & \oui & \non & \non & Auto-vérification (après) & Calcul & \non \\
Fu et al. \citeyearpar{fu2026inflation} & \oui & \non & \non & Entropie de branchement (avant, locale) & Tokens réels & \partiel \\
RACER \citeyearpar{hao2026racer} & \oui & \non & \oui & Score calibré & Taille d'ensemble & \non \\
\midrule
Mallen et al. \citeyearpar{mallen2022trust} & \non & \oui & \non & Popularité de l'entité & Inférence & \non \\
Adaptive-RAG \citeyearpar{jeong2024adaptiverag} & \non & \oui & \non & Complexité (classifieur) & Temps, étapes & \non \\
Corrective RAG \citeyearpar{yan2024crag} & \non & \oui & \non & Qualité des documents & -- & \non \\
Moskvoretskii et al. \citeyearpar{moskvoretskii2025adaptive} & \non & \oui & \non & Incertitude (27 méthodes) & Appels, récupérations & \non \\
SR-RAG \citeyearpar{wu2025srrag} & \non & \oui & \non & Incertitude + mémoire de politiques & Récupérations, latence & \partiel \\
RAGRouter-Bench \citeyearpar{wang2026ragrouterbench} & \partiel & \oui & \non & Compatibilité requête--corpus & Ressources & \non \\
\midrule
Kamath et al. \citeyearpar{kamath2020selective} & \non & \non & \oui & Calibrateur appris & -- & \non \\
Samsi et al. \citeyearpar{samsi2023watts} & \non & \non & \non & (mesure seulement) & Énergie GPU & \non \\
\midrule
\textbf{Projet visé} & \oui & \oui & \oui & \textbf{Signaux peu coûteux, coût du signal compté} & \textbf{Latence, mémoire, énergie} & \oui \\
\bottomrule
\end{tabularx}
\caption{Positionnement des travaux les plus proches. \oui{} : traité ; \non{} : non traité ;
\partiel{} : partiellement (par exemple des modèles 7B exécutables localement, sans évaluation en
contexte local). Grille établie à partir des résumés, à confirmer à la lecture.}
\label{tab:comparaison}
\end{table}

%======================================================================
\section{Synthèse et positionnement}\label{sec:positionnement}
%======================================================================

\subsection{Limites communes des travaux existants}
\begin{enumerate}
\item \textbf{Des décisions isolées.} Le routage choisit le modèle, le RAG adaptatif choisit la
récupération et l'abstention est étudiée à part. Aucun des travaux recensés ne choisit
\emph{conjointement} les trois dans une même politique (tableau~\ref{tab:comparaison}).
\item \textbf{Un coût surtout financier.} Le coût est surtout exprimé en prix d'API ou en nombre
d'appels ; la latence, la mémoire et l'énergie mesurées sur machine locale sont rarement rapportées,
une limite que souligne aussi \citet{varangot2025survey}.
\item \textbf{Le coût du signal est ignoré.} Les signaux d'incertitude les plus performants exigent
plusieurs générations, et leur coût est rarement retranché du gain.
\item \textbf{Une évaluation parfois optimiste.} Certains gains disparaissent face à un contrôle
« requête seule » \citep{li2026beyond}, et une moyenne globale peut masquer des dégradations sur les
questions difficiles.
\end{enumerate}

\subsection{Contribution potentielle}
\begin{center}
\fcolorbox{teal!60!black}{teal!5}{\parbox{0.92\textwidth}{
Concevoir et évaluer une \textbf{politique de décision conjointe et interprétable} qui choisit, pour
chaque requête, \textbf{la taille du modèle}, \textbf{le niveau de récupération documentaire} (aucune,
simple, complémentaire) et \textbf{la décision de répondre ou de s'abstenir}. Cette politique s'appuie
sur des signaux peu coûteux \textbf{dont le coût est lui-même comptabilisé}, et elle est évaluée
\textbf{sur du matériel local} en qualité, factualité, couverture/risque, latence, mémoire et, si
mesurable, énergie.}}
\end{center}

\subsection{Hypothèses de travail}
\begin{description}[style=nextline,leftmargin=0.5cm]
\item[H1 (QR1)] Un signal de difficulté calculé avant génération et à faible coût suffit à orienter une
part significative des requêtes vers le petit modèle, sans RAG, sans perte de qualité au-delà de la
marge fixée.
\item[H2 (QR2)] Combiner un signal d'incertitude du modèle et un score de suffisance des preuves
améliore le compromis couverture/risque par rapport à chacun des deux signaux pris séparément.
\item[H3 (QR3)] La politique conjointe domine les stratégies fixes sur la courbe qualité--coût
(latence, tokens et mémoire), une fois le coût de ses propres signaux inclus.
\item[H4 (QR4)] Les gains varient selon la catégorie de question ; ils seront rapportés par catégorie,
pas seulement en moyenne.
\end{description}

\subsection{Travaux dont il faut se distinguer explicitement}
\begin{center}\small
\begin{tabularx}{\textwidth}{l Y Y}
\toprule
\textbf{Travail} & \textbf{Ce qu'il fait déjà} & \textbf{Ce qui reste ouvert} \\
\midrule
Adaptive-RAG & Complexité $\rightarrow$ stratégie de RAG & Choix du modèle, abstention, mesure locale \\
Moskvoretskii et al. & Signaux d'incertitude simples pour décider du RAG & Plusieurs tailles de modèle, abstention, énergie \\
RAGRouter-Bench & Routage entre paradigmes de RAG avec coût en ressources & Abstention, matériel local, politique conjointe modèle $\times$ RAG \\
Fu et al. (2026) & Signal local avant exécution, coût réel & RAG adaptatif, abstention \\
RACER & Routage avec abstention calibrée & RAG, coût local \\
\bottomrule
\end{tabularx}
\end{center}

\subsection{Conséquences pour les phases suivantes}
\begin{itemize}
\item \textbf{Benchmark (Phase 2).} Reprendre SQuAD/SQuAD 2.0, HotpotQA et MuSiQue pour pouvoir se
comparer à Adaptive-RAG ; étudier l'intérêt de RAGRouter-Bench.
\item \textbf{Baselines (Phase 3).} En plus des baselines A, B et C du cahier des charges, ajouter
Adaptive-RAG (code public), un routeur fondé sur la seule requête (contrôle de \citealp{li2026beyond}),
et si possible un seuil d'incertitude simple (dans l'esprit de \citealp{moskvoretskii2025adaptive}).
\item \textbf{Métriques (Phase 2).} Courbes qualité--coût ; couverture/risque pour l'abstention ;
résultats par catégorie ; coût des signaux inclus dans le coût total.
\end{itemize}

%======================================================================
\section{Limites de cette revue et prochaines étapes}\label{sec:limites}
%======================================================================

\begin{itemize}
\item \textbf{Pré-analyse.} L'analyse repose sur les résumés ; les cellules du
tableau~\ref{tab:comparaison} et les limites attribuées aux travaux doivent être confirmées par la
lecture complète.
\item \textbf{Couverture.} La recherche croisée a porté surtout sur arXiv ; elle doit être complétée par
ACL Anthology, Semantic Scholar et le \textit{snowballing} à partir de RAGRouter-Bench, SR-RAG et RACER.
\item \textbf{Prépublications.} Plusieurs travaux de 2025--2026 ne sont pas encore évalués par les pairs
et peuvent évoluer.
\item \textbf{Un domaine très actif.} Une veille (alertes de citation) est nécessaire jusqu'à la
soumission ; le positionnement devra être revérifié juste avant.
\end{itemize}

\paragraph{Prochaines étapes.} (1) Lecture approfondie des travaux P1 et des cinq travaux du tableau de
distinction ; (2) mise à jour de cette revue (v0.2) ; (3) consolidation de la formulation de la
contribution, idéalement après relecture par un enseignant-chercheur ; (4) passage à la Phase 2 (benchmark et protocole).

\appendix
\section{Requêtes de recherche croisée (arXiv, 17/10/2026)}\label{ann:requetes}
\begin{center}\small
\begin{tabularx}{\textwidth}{Y c}
\toprule
\textbf{Requête (champ résumé / titre)} & \textbf{Résultats examinés} \\
\midrule
\texttt{abs:routing AND abs:retrieval AND abs:small AND abs:large AND abs:LLM} & 12 \\
\texttt{abs:"adaptive retrieval" AND abs:"model selection"} & 0 \\
\texttt{ti:RAG AND ti:rout*} & 12 \\
\texttt{abs:"retrieval-aware" AND abs:routing} & 1 \\
\texttt{abs:abstention AND abs:routing AND abs:LLM} & 12 \\
\texttt{abs:energy AND abs:routing AND abs:"retrieval-augmented"} & 9 \\
\midrule
\textbf{Total examiné sur titre / retenus après résumé} & \textbf{46 / 8} \\
\bottomrule
\end{tabularx}
\end{center}

\bibliographystyle{plainnat}
\bibliography{references}

\end{document}
